\section{Introduction}
\label{sec:intro}

% --- Background: ASR landscape ---
End-to-end speech recognition has matured from hybrid HMM-DNN systems
\cite{hinton2012deep} to large-scale weakly-supervised models such as
Whisper \cite{radford2023whisper}, which demonstrate that scaling data
and model size delivers strong recognition across languages and acoustic
conditions. In parallel, large language models (LLMs) have become the
dominant backbone for text understanding and generation, and several
recent systems -- Qwen-Audio \cite{chu2023qwenaudio},
SALMONN \cite{tang2024salmonn}, Step-Audio~2 \cite{stepaudio22025},
Kimi-Audio \cite{kimiaudio2025} and Qwen3-ASR \cite{qwen2025qwen3asr}
-- show that connecting a strong audio encoder to a pre-trained LLM
yields a single model that transcribes, understands and follows
instructions in one pass.

% --- Motivation: real-world failure modes the LLM addresses ---
The LLM backbone is more than a stronger language model: its prompt
interface is a natural place to express the kind of side information
real deployments require, and its training stack is a natural place
to fold in the acoustic robustness real deployments demand. Three
failure modes dominate the residual error of an otherwise-strong
multilingual ASR model:

\begin{itemize}
  \item \textbf{Rare named entities.} Names, jargon, product SKUs and
        long-tail terms that a generic ASR model is unlikely to have
        memorised. Classical solutions involve external biasing
        graphs or on-the-fly LM rescoring; an LLM can simply accept a
        list of likely terms as text in its prompt.
  \item \textbf{Overlapping speakers.} In multi-talker audio the user
        often wants the transcript of \emph{one} known speaker. Source
        separation followed by ASR is brittle; an end-to-end LLM that
        receives a short reference utterance can be trained to decode
        only the matching speaker.
  \item \textbf{Noisy and far-field acoustic conditions.}
        Reverberation, background noise and competing voices are
        pervasive in real recordings but rare in clean read-speech
        training sets. Augmentation curricula are the standard
        remedy; the released AmphionASR checkpoint inherits noise
        / RIR / overlap augmentation as a by-product of the
        target-speaker training pipeline, and additionally
        completes a noise-robustness SFT stage that extends the
        same augmentation to the plain-ASR branch.
\end{itemize}

% --- This work ---
We describe \textbf{AmphionASR}, an open-source, industrial-grade
SpeechLLM whose LLM prompt is the single interface for personalised
Mandarin and English speech recognition: hotword lists,
target-speaker reference audio, and noisy or overlapping-speaker
inputs are all consumed as prompt features, while the model remains
competitive on standard Mandarin and English ASR benchmarks. The
released AmphionASR checkpoint totals $4.7$\,B parameters
($4.3$\,B distinct after input--output embedding tying).

Our contributions are:

\begin{itemize}
  \item \textbf{An industrial-grade Mandarin--English SpeechLLM
        recipe.} We couple a Whisper-style 24-layer audio encoder
        ($\approx\!314$\,M parameters, initialised from the
        Qwen3-ASR audio tower) with the Qwen3-4B-Instruct
        \cite{qwen2025qwen3} language model through a lightweight
        $\approx\!9.2$\,M-parameter modality projector (a two-layer
        MLP with the SwooshR activation \cite{yao2024zipformer})
        and four learnable modality-boundary prompts. On a curated
        public-plus-internal corpus the model reaches $1.85$\,\%
        CER on AISHELL-1, $1.82$\,\% WER on LibriSpeech
        \texttt{test-clean} and $7.48$\,\% MER on the TALCS
        Mandarin--English code-switching test set
        \cite{li2022talcs} (Section~\ref{sec:experiments}).
  \item \textbf{Strong oracle-prompt hotword conditioning.} With
        five in-prompt hotwords the CommonVoice~17 hotword test
        CER drops from $6.01$\,\% to $2.67$\,\% on Mandarin and
        the WER from $7.68$\,\% to $4.69$\,\% on English (oracle
        protocol, random padding;
        Section~\ref{sec:hotwords}).
  \item \textbf{A parameter-free hotword retrieval operator.} We
        further describe a text-space retrieve-then-bias operator
        that, paired with a matched SFT hard-negative augmentation,
        scales prompt-driven hotword behaviour to pools of size up
        to $10^4$. The retriever attains Recall@$K$ above $94\%$
        on Mandarin from $K=10$ onwards and Recall@$K \approx 78\%$
        on English at $K=20$
        (Section~\ref{subsec:retrieve-algorithm}).
  \item \textbf{Target-speaker recognition with reference audio.}
        A $3$--$5$\,s reference utterance suffices to make the
        model attend to a single speaker in mixed audio. On our
        internal hot+TS test set we reach $13.06$\,\% MER, with
        explicit anti-hallucination training that keeps the
        decoder silent when the requested speaker is absent; we
        report and analyse the substantially weaker numbers on the
        public LibriMix benchmarks (Section~\ref{sec:tsasr}).
  \item \textbf{A noise-robustness SFT stage.} The TS-ASR branch
        already passes through reverberation, additive noise
        (MUSAN \cite{snyder2015musan} and AudioSet road-traffic)
        and competing-speaker interference; an additional
        noise-robustness SFT stage extends the same augmentation
        to the plain-ASR branch and gives AmphionASR its first
        controlled noise-robustness evaluation grid against
        Qwen3-ASR-1.7B~\cite{qwen2025qwen3asr}
        (Section~\ref{sec:noise}).
\end{itemize}

% --- Roadmap ---
Section~\ref{sec:related} surveys related work; Section~\ref{sec:data}
documents the training corpus; Section~\ref{sec:model} introduces the
architecture; Section~\ref{sec:training} details the staged training
recipe; Section~\ref{sec:experiments} reports the multilingual main
results; Sections~\ref{sec:hotwords}--\ref{sec:noise} dig into the three
flagship capabilities; Section~\ref{sec:analysis} consolidates failure
modes and ablations; we discuss limitations in
Section~\ref{sec:limitations} and conclude in
Section~\ref{sec:conclusion}.
